% Source: source/普通高考/2008/2008江苏.pdf, page 1, question 7.
% Nodes: 开始; S←0; i←1; 输入 G_i,F_i; S←S+G_i×F_i; i≥5;
% 输出 S; 结束; i←i+1.
% Directed edges: 开始→S←0→i←1→输入→累加→判定;
% 判定(Y)→输出 S→结束; 判定(N)→i←i+1→输入.
\begin{tikzpicture}[
  >=Stealth,
  line width=0.55pt,
  line cap=round,
  line join=round,
  every node/.style={font=\normalsize,anchor=center,align=center,
    execute at begin node={\everypar{}\leftskip=0pt\rightskip=0pt\parindent=0pt
      \hangindent=0pt\hangafter=1\parshape=0}},
  flowbox/.style={draw,minimum height=.70cm,inner xsep=4pt,inner ysep=2pt,
    align=center},
  terminal/.style={flowbox,rounded corners=.18cm,minimum width=1.55cm,
    minimum height=.74cm},
  io/.style={flowbox,shape=trapezium,trapezium left angle=72,
    trapezium right angle=108,minimum height=.80cm},
  decision/.style={flowbox,shape=diamond,aspect=2.85,minimum width=3.45cm,
    minimum height=1.00cm,inner sep=1pt},
  flowarrow/.style={->,line width=0.55pt},
]
  \node[terminal] (start) at (0,10.90) {开始};
  \node[flowbox,minimum width=2.25cm] (initS) at (0,9.55)
    {$S\leftarrow 0$};
  \node[flowbox,minimum width=2.15cm] (initI) at (0,8.20)
    {$i\leftarrow 1$};
  \node[io,minimum width=3.85cm] (input) at (0,6.80)
    {输入 $G_i,F_i$};
  \node[flowbox,minimum width=4.60cm] (accum) at (0,5.30)
    {$S\leftarrow S+G_i\times F_i$};
  \node[decision] (test) at (0,3.72) {$i\geqslant 5$};
  \node[flowbox,minimum width=2.70cm] (increment) at (-4.05,5.30)
    {$i\leftarrow i+1$};
  \node[io,minimum width=2.35cm] (output) at (0,2.02) {输出 $S$};
  \node[terminal] (finish) at (0,0.68) {结束};

  \draw[flowarrow] (start.south) -- (initS.north);
  \draw[flowarrow] (initS.south) -- (initI.north);
  \draw[flowarrow] (initI.south) -- (input.north);
  \draw[flowarrow] (input.south) -- (accum.north);
  \draw[flowarrow] (accum.south) -- (test.north);

  % 否分支回到输入框；回线不与输出分支共线。
  \draw[flowarrow] (test.west) -- (-4.05,3.72) -- (increment.south);
  \draw[flowarrow] (increment.north) -- (-4.05,6.80) -- (input.west);
  % Keep the N branch label clear of the horizontal connector and diamond edge.
  \node[anchor=east] at (-1.70,4.04) {N};

  % 是分支独立通往输出，不与循环回线交汇。
  \draw[flowarrow] (test.south) -- (output.north);
  \node[anchor=west] at (0.18,2.86) {Y};
  \draw[flowarrow] (output.south) -- (finish.north);

  \path[use as bounding box] (-5.45,0.30) rectangle (2.35,11.30);
\end{tikzpicture}
